% ============================================================
% Section 1.1: General Overview and Clinical Background
% ============================================================
\section{General Overview and Clinical Background}
\label{sec:general_overview_clinical_background}

Severe neurodegenerative disorders and traumatic injuries---such as advanced Amyotrophic Lateral Sclerosis (ALS), Muscular Dystrophy, brainstem stroke, and high-level spinal cord injuries---frequently result in profound motor paralysis, often presenting as Locked-in Syndrome (LIS) or Quadriplegia. Patients in these clinical conditions suffer from a complete loss of voluntary muscle control throughout the body, depriving them of conventional speech and physical interaction with their environment. However, the cranial nerves controlling extraocular muscle motility remain physiologically intact and functionally operational in most affected individuals.

To address these severe functional deficits, non-invasive Human-Machine Interfaces (HMIs) based on electrooculography (EOG) have emerged as an effective assistive solution. EOG exploits the standing electrostatic potential maintained between the cornea and the retina to track eye displacement and voluntary eyelid blinks.

